\section{Related Work}
\label{sec:related_work}

\subsection{Dermoscopic Classification and Backbones}

Public ISIC resources and HAM10000 have supported reproducible research on
dermoscopic classification
\cite{codella2019isic2018,tschandl2018ham10000}. EfficientNet jointly scales
network depth, width, and resolution and provides compact baseline
architectures \cite{tan2019efficientnet}. DenseNet connects each layer to
preceding feature maps, encouraging feature reuse and information flow
\cite{huang2017}. In this work, EfficientNet-B0 provides the backbone-matched
flat-versus-hierarchy comparison, while DenseNet-121 tests whether the
metric-specific findings persist under a protocol-matched alternative
backbone.

\subsection{Class Imbalance}

Focal loss reduces the influence of well-classified examples and emphasises
difficult observations \cite{lin2017focal}. Class-balanced loss derives
weights from the effective number of samples and is designed for long-tailed
data \cite{cui2019classbalanced}. Such methods remain insufficiently
characterised by accuracy alone; macro-F1, balanced accuracy, confusion
matrices, and class-specific recall are necessary to expose minority-class
behaviour.

\subsection{Hierarchical Classification}

Top-down hierarchical classifiers are vulnerable to error propagation because
an incorrect decision at an upper level constrains all subsequent predictions
along the selected path \cite{silla2011survey}.
Hsu and Tseng incorporated diagnostic hierarchy information through
hierarchy-aware contrastive learning \cite{hsu2022hierarchy}. More recently, Hsu and Tseng proposed LightDPH, a lightweight
hierarchical contrastive-learning framework that jointly considers broad
lesion groups and fine-grained classes without relying on a conditional
inference gate \cite{hsu2024lightdph}. Suleiman
\emph{et al.} evaluated a two-step hierarchical skin-lesion system combining
transfer learning and conventional classifiers
\cite{suleiman2024twostep}. These approaches demonstrate the potential of
diagnostic structure, but representation-level hierarchy and conditional
routing are distinct. In a conditional system, an upstream decision determines
whether a downstream classifier is executed, creating an explicit
error-propagation path.

The present study therefore evaluates both endpoint performance and the
routing mechanism. Predicted routing is compared with oracle routing while the
trained downstream classifier remains fixed, allowing routing-associated loss
to be separated from subtype errors after correct routing.

\subsection{Paired Statistical Evaluation}

Predictions from models evaluated on the same images are paired. Paired
bootstrap resampling can estimate uncertainty in metric differences
\cite{efron1993bootstrap}, while McNemar's test evaluates disagreement in
paired correctness outcomes \cite{mcnemar1947sampling}. These procedures were applied to all pairwise model comparisons on the
common locked test cohort, with metric-specific interpretation of the
resulting confidence intervals and paired correctness tests.
